% This is a template for a thesis (bachelor or master)
% from CogSys Group of Otto-Friedrich-Universität Bamberg.
%
% Copyright: CogSys Group
% Last modification: 2021-07-26
%
\documentclass[
%english,% if you're writing in English (default),
%german,% if you're writing in German
thesis% please do not delete or change this line!
]{csthes}

% The following packages are loaded automatically. Do not load them
% manually as you might risk an option clash:
%   -babel (heading translation and hyphenation)
%   -natbib (citing, see www.ctan.org/tex-archive/macros/latex/contrib/natbib/natbib.pdf)
%   -fancyhdr (headers and footers)
%   -fontenc (output encoding)

% Set the encoding your tex-source is written in.
% Use utf8 if possible. Other options that might 
% work for you are:
%   latin1 (Unix/Linux oder Windows)
%   ansinew (Windows)
%   applemac (Mac)
\usepackage[utf8]{inputenc}

%for tables
\usepackage[table,xcdraw]{xcolor}

%for graphics
\usepackage{graphicx}

% For the online version enable hyperlinks
\usepackage{hyperref}

%%%%%%%%%%%%%%%%%%%%%%%%%%%%%%%%%%%%%%%%%%%%%
%  Set informations about your thesis here  %
%%%%%%%%%%%%%%%%%%%%%%%%%%%%%%%%%%%%%%%%%%%%%
\title{Your Title, e.g. Rule Learning for LaTeX Templates}
\subtitle{Your Subtitle, e.g. Comparing Different Learning Strategies}% optional
\setthesistype{Bachelor/Master}
\author{Your Name, e.g. Zaphod Beeblebrox \matr{Your Matriculation Number}}

\degree{Your Course of Study, e.g. Applied Computer Science}
\supervisor{Prof.\,Dr.\,Ute Schmid}% most likely

% Start of content
\begin{document}
% First set the title page
\maketitle
% Give a short summary of your thesis
\begin{abstract}
A summary of your work, about half a page of length. An abstract should give the reader all relevant information about the paper in a nutshell: Topic, research question, method, results, evaluation of results.
\end{abstract}
% Set the table of contents
\tableofcontents%

% And other tables/lists optionally
\listoftables%
\listoffigures%

% Start the main work (handles page numbering etc.)
\mainmatter

%%%%%%%%%%%%%%%%%%%%%%%%%%%%%%%%%%%%%%
%  Place your content here           %
%                                    %
%  or even better:                   %
%  use \input{chapters/chaptername}  %
%  for each chapter!                 %
%%%%%%%%%%%%%%%%%%%%%%%%%%%%%%%%%%%%%%

% While writing your thesis it is advisable to only load the chapter you are currently working on. It decreases the compile time. 

% 1. Chapter: ''Introduction''
\input{chapters/first_chapter}
% 2. Chapter: ''State of Research''
\input{chapters/second_chapter}
% 3. Chapter: ''Methods''
\input{chapters/third_chapter}
% 4. Chapter: ''Realisation and Evaluation''
\input{chapters/fourth_chapter}
% 5. Chapter: ''Conclusions and Future Work''
\input{chapters/five_chapter}

% The chapters Introduction and Conclusions have NO sections
% For the other chapters:
% Each sectioning (chapter, section, subsection, ...) should have NONE
% or AT LEAST TWO sub parts!
%\section{First Section}
%\subsection{First Subsection}
%\subsubsection{First Subsubsection}
%\paragraph{First Paragraph.} Some content for the first paragraph. 
%\subparagraph{First Subparagraph.} Some content for the first subparagraph. 
% Please note the dot after the paragraph and subparagraph heading. This is not accidently.

%Use bibtex for references to literature. 
%The most references will be in the second chapter.

Some further recommendations: use diagrams, images and tables to illustrate your solution and results. If you do so, please add captions with citations whenever a resource was not created by you. For tables, it is a good style to not use horizontal lines. Some example table:

% Please add the following required packages to your document preamble:
% \usepackage[table,xcdraw]{xcolor}
\begin{table}[hb!]
\centering
\caption{Table according to (add citation)} 
\begin{tabular}{|c|l|l|l|l|}
\hline
\rowcolor[HTML]{C0C0C0} 
\multicolumn{1}{|l|}{\cellcolor[HTML]{C0C0C0}{\color[HTML]{000000} Some Table}} & {\color[HTML]{000000} A} & {\color[HTML]{000000} B} & {\color[HTML]{000000} C} & {\color[HTML]{000000} D} \\ 
A                                                                               & x                        &                          &                          &                          \\ 
B                                                                               &                          &                          & x                        &                          \\ 
C                                                                               &                          & x                        &                          & x                        \\ 
\hline
\end{tabular}

\end{table}

Some example image:
\begin{figure}[hb!]
\centering
  \includegraphics[width=0.4\textwidth]{images/UB-Logo-blau}
    \caption{The logo of your university}
\end{figure}

% Make the references section
\bibliography{mybib} % According to the name of your bib file

% Start the appendix
\appendix
%%%%%%%%%%%%%%%%%%%%%%%%%%%%%%%%%%%%%%%%%%%
%  Place appendix content here            %
%  or use \input{appendix/appendixname}   %
%%%%%%%%%%%%%%%%%%%%%%%%%%%%%%%%%%%%%%%%%%%
\chapter{First Appendix Chapter}
\section{First Appendix Section}
\subsection{First Appendix Subsection}
\subsubsection{First Appendix Subsubsection}

\cleardoublepage

%----------------------------------------------------------------------------------------
%	DECLARATION PAGE
%---------------------------------------------------------------------------------------

% TODO uncomment the next three lines if your thesis is wirtten in german 
% \chapter*{Eidesstattliche Erklärung}
% \markboth{Eidesstattliche Erklärung}{Eidesstattliche Erklärung}
% \addcontentsline{toc}{chapter}{Eidesstattliche Erklärung}

% TODO comment the next three lines if your thesis is written in german 
\chapter*{Declaration of Authorship}
\markboth{Declaration of Authorship}{Declaration of Authorship}
\addcontentsline{toc}{chapter}{Declaration of Authorship}
% TODO Change the declaration according as needed. 

Ich erkläre hiermit gemä\ss\ \S~9 Abs.\,12 APO, dass ich die vorstehende {\thesistype}arbeit 
selbständig verfasst und keine anderen als die angegebenen Quellen und Hilfsmittel benutzt habe. 
Des Weiteren erkläre ich, dass die digitale Fassung der gedruckten Ausfertigung 
der {\thesistype}arbeit ausnahmslos in Inhalt und Wortlaut entspricht und zur Kenntnis 
genommen wurde, dass diese digitale Fassung einer durch Software unterstützten, 
anonymisierten Prüfung auf Plagiate unterzogen werden kann. 

\bigskip
\bigskip

\begin{tabular}{@{}l@{}}
  Bamberg, den \rule[-0.8em]{7em}{0.5pt}\\[2ex]
  ~
\end{tabular}
\hspace{\fill}%
\begin{tabular}{@{}c@{}}
  \rule[-0.8em]{19em}{0.5pt}\\[2ex]
  {Your name e.g. Walentina Tereschkowa-Nikolajewa}
\end{tabular}\hspace{\fill}

% Finish content and document
\end{document}
